% !TeX root = ../../../../../main.tex
% chapters/sections/chapter4/subsections/method/occbin.tex

\subsection{Occbinの役割}
\label{sec:method_occbin}
前述の陰関数\(h\)は、方程式系に含まれる名目利子率などの変数が制約なく自由に変動できるという前提の下で
導出されたものである。
しかし本モデルにおいては、名目利子率の下限が0に制約されるゼロ金利制約を想定している。
通常時において名目利子率はテイラールールなどの金融政策の反応関数に従って決定されるが、
名目利子率が0を下回る状態に陥ると実際の金利は0に固定される。
これは数学的に見れば、連立方程式系を構成する1つの式（金融政策の反応関数）が、
「\(i_t=0\)」という独立した式に置き換わることを意味する。
式が置き換わることで、インフレ率や生産量などに対する金利の反応係数がすべて0となるため、
前節で説明した残差関数\(f\)の偏微分から構成される係数行列\(A,B,C,D\)において、
当該方程式に該当する行の値が変化する。
係数行列の値が変われば、そこから代数的に導出される推移行列\(W\)の値やその固有値も変化するため、
経済の動学的な性質が根本的に変わってしまう。
したがって、変数が自由であることを前提として導出した単一の陰関数\(h\)を用いて、
ゼロ金利制約下の動学を描写することは不可能となる。
そこでOccbinは、このように適用される方程式の組み合わせが異なる状態をそれぞれ「レジーム」として定義する。
具体的には、名目利子率が自由に変動する「通常レジーム」と、
名目利子率が\(i_t=0\)に固定される「制約レジーム」を設ける。
そして、それぞれのレジームにおける方程式系に対して前節と同様に1次近似をおこない、
通常レジーム用の陰関数\(h\)とは推移行列の異なる制約レジーム用の陰関数を個別に導出する。
シミュレーションにおいて、Occbinは「予測と検証」と呼ばれるアルゴリズムを用いてこれらを繋ぎ合わせる。
具体的には、まずショック発生後から定常状態に戻るまでのすべての期間において
通常レジームの陰関数\(h\)が適用されると仮定して、変数の将来経路を計算する。
そして、その経路の中で名目利子率の予測値が0を下回る期間を発見した場合、その期間を制約レジームに割り当てる。
次に、その割り当てに従って、制約期間中は制約レジーム用の陰関数を、
そこから抜け出した後の期間は通常レジーム用の陰関数\(h\)を適用して将来経路を再計算する。
再計算された経路においてレジームの仮定と金利の実現値に矛盾がないかを確認し、
矛盾がなくなるまでレジーム期間の予測と経路の再計算を繰り返すことで、
ゼロ金利制約という非線形な現象のシミュレーションを可能にしている。